\section{Introduction}
\label{sec:intro}

The United States Congress consists of 435 House members distributed among
50 states. State legislatures collectively contain approximately 7,400 seats
in 99 chambers---roughly 17$\times$ as many elected positions, drawing maps
for districts that are 17$\times$ smaller and correspondingly more sensitive
to local geography. Despite this greater scale and civic significance, state
legislative redistricting has received far less systematic attention from the
algorithmic redistricting community.

This paper presents the first comprehensive 50-state empirical study of
algorithmically generated state house maps. We apply the minimum-edge-cut
recursive bisection algorithm implemented in the \texttt{redist} Rust binary
\citep{dellaLibera2026recursive} to all 50 lower legislative chambers,
adapting the resolution (census tract vs.\ block group) based on the
$k/n_{\mathrm{tracts}} > 0.05$ threshold derived in companion paper F.3
\citep{dellaLibera2026stateLegResolution}.

\paragraph{Main contributions.}
\begin{enumerate}
  \item A 50-state database of algorithmically generated state house maps
    at the 2020 census cycle, freely available for replication.
  \item Systematic compactness measurement showing $\overline{PP} = 0.401$
    at block-group resolution---exceeding the congressional baseline of
    $\overline{PP} = 0.361$ from \citet{dellaLibera2026mec}.
  \item Population balance within $\pm 0.5\%$ for all 50 states, achieved
    despite the \emph{Reynolds v.\ Sims} standard requiring only $\pm 5\%$.
  \item County-split analysis showing 15--25\% fewer splits than enacted maps
    for the 38 states with available enacted 2020 state house maps.
  \item Four case studies (WA, TX, NH, NE) examining how district count,
    geography, and resolution interact to shape algorithmic map outcomes.
\end{enumerate}

\paragraph{Scope.}
This paper covers lower chambers (state houses and assemblies) only.
Upper chambers (state senates) and the bicameral nesting problem are
addressed in companion paper F.2 \citep{dellaLibera2026bicameral}.
VRA compliance in state legislative redistricting is addressed in F.6
\citep{dellaLibera2026vraStateLeg}.

\paragraph{Organization.}
Section~\ref{sec:related} reviews the state legislative redistricting
literature. Section~\ref{sec:methodology} describes the methodology.
Section~\ref{sec:results} presents results. Section~\ref{sec:cases}
develops four case studies. Section~\ref{sec:comparison} compares to
enacted maps. Section~\ref{sec:limitations} discusses limitations.
Section~\ref{sec:conclusion} concludes.
